Da Hesel modellen består af et koblet system af partielle differentialligninger, som ikke har en analytisk løsning, approksimeres løsningen numerisk.
Fremgangsmåden er at bruge fysikken bag Hesel modellen og data simuleringer af dynamikkerne til at træne en model $\hat{f}_{\theta}$, til at minimere dynamikker og PDE fejl.
Det forvæntes at en model trænet på et bestemt sæt af dynamikker og PDE, ville kunne genskabe både interpolation og ekstrapolation simuleret dynamikker, og samtidig opfylde PDE systemet.
Hvis ekstrapolation ikke er så godt vil modellen kunne trænes samtidig med simuleringen køre og kunne blive brugt meget hurtigt til dette nye sæt af dynamikker og PDE system.
Ved brug af disse metoder forvæntes det at simuleringen kan gøres computationalt billigere, hurtigere og samtidig vedligeholde fysisk nøjagtighed.

Den videnskabelige tilgang er som følgene:
\begin{itemize}
    \item Et grundigt litteraturstudie af PDE'er, numeriske metoder, datadrevne og fysik informede metoder
    \item Forstå og redgøre for de grundlæggende matematiske og fysiske principper bag Hesel modellen
    \item Udsvikle ablations studier og eksperimenter
    \item Implementere og træne datadrevne og fysik informede modeller
    \item Evaluere og sammenligne resultaterne
    \item Iterere og forbedre modellerne baseret på resultaterne
\end{itemize}
Disse approksimative modeller bliver evalueret på følgende test:
\begin{itemize}
    \item Hvor godt de approksimative modeller opfylder Hesel modellen
    \item Genskabe både interpolation og ekstrapolation simuleret dynamikker
\end{itemize}